\cotefile{156-4}
\scanpage{58}\mspage{57}
\diagram{\begin{tikzcd}[ampersand replacement=\&,column sep=large]
\underset{\substack{\text{ss-fig de}\\ F}}{\mathcal F_{\leq F}}
 \times
\underset{\substack{\text{ss fig}\\ \text{de }G}}{\mathcal F_{\leq G}}
 \arrow[r,"\sim"] \&
\underset{\substack{\text{ss. fig}\\ \text{de }F\cup G=H}}{\mathcal F_{\leq F\cup G}}
\end{tikzcd}}

(On récupère $F',G'$ : partir de $H'=F'\cup G'$ comme $H\cap F$ et $H\cap G$.) It. --- pour \unsure{sup.|une} \deleted{\unsure{familles|famille} mutuellement disjointes d'une} \inserted{\inserted{des} figures \unsure{pleines} : \unsure{donc} des parties \unsure{f.|fermées} $\unsure{\searrow|x}$} configuration (famille pas \unsure{néc.t|néct} finie). \deleted{\unsure{figure|figures} (\unsure{y}}

\uuline{Démonstration} \unsure{Particulièrt|Particulièrmt} \unsure{claire|clair}, quand on interprète $\mathcal F$ en termes de l'ens $\Phi\subseteq\mathcal P(M)$ \unsure{satisfais.} $\to$ a) b) c). L'ens ordonné $M_{\leq F\vee G}$ s'identifie à la réunion disjointe dans l'ens ordonné somme de $M_{\leq F}$ et $M_{\leq G}$, \unsure{donc|de} ses parties fermées \deleted{\unsure{sont|de $M$}} \unsure{s'en|si on} \unsure{définies|exprime} pr les couples de parties fermées de $M_{\leq F}$ et $M_{\leq G}$ \ldots

\uline{Remarque}
% slanted
\margin{à joindre à la remarque de la page 55}
En termes de l'ens \uline{ordonné} $M$, la donnée de $\Phi$ équivaut : un ens de parties \uline{fermées} de $M$, satisfaisant
\[
\begin{aligned}
\text{a) }&\deleted{\Phi}\forall\,\deleted{\unsure{=}}A\in\Phi,\quad
 \deleted{\text{\unsure{toute}}}\inserted{\text{et}}\ \text{partie fermée }A'\text{ de }A,\ \text{on a }A'\in\Phi\\
&\deleted{\text{b) }\phi\in\Phi\quad
 (\Longleftrightarrow\Phi\ne\phi,\ \text{compte tenu de a), et}}\\
\text{b) }&\bigcup_{A\in\Phi} A=M\quad\deleted{\unsure{\forall x}}\quad
 \left[\,\underset{\substack{\text{à cause}\\\text{de a)}}}{\Longleftrightarrow}\
 \forall x\in M,\ M_{\leq x}\in\Phi\,\right]\\
\text{c) }&\phi\in\Phi\quad
 \bigl(\underset{\deleted{\unsure{\text{a)}|\text{a) b)}}}}{\Longleftrightarrow}
 \Phi\ne\phi,\ \text{et \unsure{automatiqt} satisfaite}\ \deleted{\unsure{s|\Phi}}\\
 &\hspace{7em}\text{(à cause de b)) si }M\ne\phi\bigr)
\end{aligned}
\]

\uuline{NB} Soit $\mathcal P_f(M)$ l'ens des parties fermées de $M$, $\unsure{\mathcal P_{f'}|\mathcal P'_f}(\mathcal F)$ de $\mathcal F$ qui sont stables par Sup. Si $\mathcal F'$ et $M'$ se correspondent \margin{TSVP}, 
\diagram{\fbox{\begin{tikzcd}[ampersand replacement=\&,column sep=large]
\mathcal P_f(\deleted{\unsure{M}}\mathcal F) \arrow[r,shift left=.6ex,"\varphi"] \& \mathcal P_f(M) \arrow[l,shift left=.6ex,"\psi"]
\end{tikzcd}}}
$\varphi(\mathcal F')=\mathcal F'\cap M)$, $\psi(M')=\{F\in\mathcal F\mid\overline F\unsure{\subset|\in} \unsure{M'|M}\unsure{\}}$ alors $\varphi,\psi$ bij. inverses \unsure{l'une|l'un} de l'autre
